\chapter{The covariant phase space and its boundary charges}\label{ch:05}
The endpoint term in an action variation does more than supply canonical
momenta. Its change under a second variation tells us whether a proposed
gauge transformation is a redundancy of a region or changes a measurable
flux at its boundary. No canonical coordinates on the full space of fields
are needed to make the test.

\section{From endpoints to a two-form}
At prescribed cut datum $j$, the \emph{pre-phase space}
$\widetilde{\mathcal P}_R[j]$ is the space of solutions satisfying the physical
and prescribed cut conditions, before any identification by gauge. At a
regular solution its tangent vectors solve the linearized equations and the
linearized boundary conditions. A spacetime integral can still be a form on
this space: for instance
\[
 \theta_\Sigma(Y)=\int_\Sigma\Theta(\phi;Y)
      -\sum_F\int_{\partial_F\Sigma}C_F(\phi;Y)
\]
is a number linear in the variation $Y$. The sum includes each side face of
Chapter~\ref{ch:04}. The exterior derivative in \emph{field} space is
\begin{align}
 \omega_\Sigma&=\delta\theta_\Sigma,\nonumber\\
 \omega_\Sigma(Y,Z)&=Y[\theta_\Sigma(Z)]-Z[\theta_\Sigma(Y)]
                                      -\theta_\Sigma([Y,Z]).
 \label{eq:05-omega}
\end{align}
In particular, $\theta=p\,\delta q$ gives $\omega=\delta p\wedge\delta q$,
not its negative. The commutator term allows $Y$ and $Z$ to vary from one
solution to the next. The integrals of $\delta C_F$ enter $\omega_\Sigma$
with exactly the same minus signs as the $C_F$ enter $\theta_\Sigma$.

Pull the first-variation identity back to the space of solutions. It becomes
$\delta S_R=\theta_{\Sigma_f}-\theta_{\Sigma_i}$. Taking one more field-space derivative gives
$0=\delta^2S_R=\omega_{\Sigma_f}-\omega_{\Sigma_i}$. This proof of slice
independence would have failed had we dropped a physical corner momentum.
For example, return to the string with a massive endpoint $B$ in
Chapter~\ref{ch:04}, holding the other endpoint fixed. For two perturbations
$u,v$, the bulk two-form is $\int(\dot u v-\dot v u)\,dx$. Differentiate it
in time and use $\ddot u=\partial_x^2u-m^2u$ and the same equation for $v$.
The mass terms cancel, and spatial integration by parts leaves
$(\partial_n u)v-(\partial_n v)u$ at $B$. The endpoint equation converts
this into $-d[\kappa(\dot u v-\dot v u)]/dt$. Thus the bulk form alone
changes by exactly the amount stored at the endpoint.

The same endpoint argument tells us what changes when we open the cut.
There the first variation
has an additional one-form $\langle\Pi,\delta q\rangle_\Gamma$, so
\[
 (\omega_{\Sigma_f}-\omega_{\Sigma_i})(Y,Z)
 =-\langle Y\Pi,Zq\rangle_\Gamma+\langle Z\Pi,Yq\rangle_\Gamma.
\]
Here the integrated potentials are those of the opened action. Flux through
one region need not vanish. On matched solutions the two fluxes cancel,
since both variations preserve the common trace and the vanishing sum of
outward responses. Closedness, $\delta\omega_\Sigma=0$, follows from its
definition even before one asks about slice independence. A closed two-form
with a possible kernel is called \emph{presymplectic}. Only after identifying
its null directions and, when regular, taking the appropriate quotient does
one obtain a symplectic physical phase space.

\section{An action symmetry and the sign of its generator}
Let $X$ be an infinitesimal transformation tangent to the specified domain,
with any external symmetry parameter held fixed when fields vary. Test it as
an action symmetry \emph{off shell}:
\[
 \iota_X\delta S_R=\alpha_{X,\Sigma_f}-\alpha_{X,\Sigma_i}+\beta_X,
 \qquad \delta\beta_X=0.
\]
The endpoint functional $\alpha_X$ includes contributions from a changed
boundary Lagrangian or corner. It need not be zero merely because the field
equations are invariant. Comparing this identity with the on-shell variation
gives the Noether generator \cite{HW}
\begin{equation}
 H_{X,\Sigma}=\iota_X\theta_\Sigma-\alpha_{X,\Sigma}.
 \label{eq:05-noether}
\end{equation}
It has equal endpoint values if $\beta_X=0$; a nonzero field-independent
$\beta_X$ shifts their difference by a prescribed number. To obtain a
Hamiltonian vector field, the transformation must also have a differentiable
charge with the chosen boundary data. For the representative of the potential
being used, the endpoint check is
$\mathcal L_X\theta_\Sigma=\delta\alpha_{X,\Sigma}$ on solution variations
(including the corner). Cartan's identity then fixes the slot and the sign:
\[
 \delta H_{X,\Sigma}=
 \mathcal L_X\theta_\Sigma-\iota_X\delta\theta_\Sigma-\delta\alpha_{X,\Sigma}
 =-\iota_X\omega_\Sigma,\qquad
 \delta H_{X,\Sigma}(Y)=\omega_\Sigma(Y,X).
\]
If the endpoint check leaves a nonzero one-form, the stated charge does not
generate $X$ with this potential and these boundary data. The test is local
to each slice; an identity for the difference of two endpoints is not by
itself this test.

For example, a free particle with $L=m\dot q^2/2$ has the boost $Xq=vt$.
Off shell $XL=mv\dot q=d(mvq)/dt$, so $\alpha_X=mvq$. Since
$Xp=mv$ and $\mathcal L_X(p\,\delta q)=mv\,\delta q$ at fixed $t$,
\eqref{eq:05-noether} gives $H_X=v(tp-mq)$. Directly,
$\omega(Y,X)=v(t\,Yp-m\,Yq)$. The term $-mvq$ is what generates the change
of momentum; $vtp$ alone would only generate the change of position.

\section{Gauss' law and a charged cut}
Take source-free Maxwell theory on an ultrastatic slab,
$g=-dt^2+h_{ij}dx^i dx^j$, with the spatial metric fixed. Vary
$L=-F_{\mu\nu}F^{\mu\nu}\epsilon/4$, $F=dA$. Antisymmetry gives
\[
 -F^{\mu\nu}\nabla_\mu\delta A_\nu
 =(\nabla_\mu F^{\mu\nu})\delta A_\nu
     -\nabla_\mu(F^{\mu\nu}\delta A_\nu).
\]
Use walls with no corner kinetic term here, for instance fixed tangential
$A$ on a physical wall; take gauge parameters to vanish there. Define
$E^i=F^{i0}=h^{ij}F_{0j}$, so that $-F^{0i}\delta A_i=E^i\delta A_i$ is
the \emph{final}-slice coefficient. With $dV_h$ the spatial volume measure,
\[
 \theta_\Sigma=\int_\Sigma E^i\delta A_i\,dV_h,
 \qquad \omega_\Sigma=\int_\Sigma\delta E^i\wedge\delta A_i\,dV_h,
 \qquad D_iE^i=0.
\]
The last equation is the time component of Maxwell's equations; $D_i$ is
the connection of $h$. The gauge transformation $X_\lambda A=d\lambda$
leaves $F$, and hence $\theta_\Sigma$, invariant for fixed $\lambda$;
$\alpha_\lambda=0$. Integrating its generator by parts gives
\begin{align*}
 H_{\lambda,\Sigma}
 &=\int_\Sigma E^i D_i\lambda\,dV_h
 =-\int_\Sigma\lambda D_iE^i\,dV_h
    +\int_{\partial\Sigma}\lambda E_\perp\,dS,\\
 \omega_\Sigma(Y,X_\lambda)
 &=\int_{\partial\Sigma}\lambda\,Y E_\perp\,dS
 \quad\text{on the linearized Gauss constraint}.
\end{align*}
Here $E_\perp=r_iE^i$, with $r$ the \emph{outward} unit normal to
$\partial\Sigma$ within $\Sigma$. For a transformation preserving the boundary
problem, vanishing trace of $\lambda$ on every boundary component is
sufficient for a proper gauge variation: it lies in the
kernel and carries no boundary charge. At a cut $S=\Gamma\cap\Sigma$ an
allowed $\lambda$ with nonzero trace can instead generate the charge
$Q_\lambda=\int_S\lambda E_\perp\,dS$ if that flux can vary. Not every
nonzero parameter gives a nontrivial generator. For example, a spatial
constant has $d_\Sigma\lambda=0$ and sees only the total flux through the
whole closed boundary, which vanishes by Gauss' law. Also, a transformation
that changes prescribed data is not tangent to that fixed-$j$ problem: it
may act between fibers in the opened family. A slice charge there need not
be conserved in the presence of flux.

For a source-free join the outward electric fluxes on the two copies of $S$
are opposite. Equal gauge parameters on both sides have zero \emph{sum} of
cut charges; the regional charge was not an error to be discarded in advance.
Nor does a surface integral by itself put independent fields at the cut. To
see what the latter choice entails, introduce a boundary coordinate $u$ and
an independent momentum density $p$, adding $\int_S p\,\delta u$ to
$\theta_\Sigma$. If $X_\lambda u=\lambda|_S$ and $X_\lambda p=0$, the new
cut generator on Gauss data is $\int_S\lambda(E_\perp dS+p)$. Imposing
$p=-E_\perp dS$ makes this diagonal transformation null on the constraint
surface. We added variables, a symmetry action, and a constraint. None of
these follows just by integrating the original bulk Gauss law by parts.

\begin{exercise}
In one spatial dimension take $h=dx^2$ on $[0,L]$ and set $A_t=0$ at the
endpoints. Set $a=\int_0^L A_x\,dx$. What charge generates time-independent
endpoint gauge parameters, and what happens if the endpoint values agree?
\end{exercise}
\noindent\emph{Answer.} Gauss' law sets $E^x=e$ independent of $x$, giving
$\theta=e\,\delta a$. The transformation shifts $a$ by
$\lambda(L)-\lambda(0)$, so $H=e[\lambda(L)-\lambda(0)]$. Equal values give
zero even if neither value vanishes. The minus sign at the left endpoint is
its outward orientation, not a sign assigned to its gauge parameter.
